% \begin{frontespizio}
% \Universita{Verona}
% \Logo{images/univr.png}
% \Facolta{Pennutologia}
% \Corso{Nanoscienze e tecnologie avanzate}
% \Dipartimento{ciccio}
% \Annoaccademico{2013-2016}
% \Titolo{Graph analysis of brain functional connectivity}
% \Sottotitolo{Breaking the resolution limit with Surprise}
% \Candidato[123456]{Carlo Nicolini}
% \Relatore{Angelo Bifone}
% \Relatore{Pasquina Marzola}
% %\Correlatore{Ugo Frogio}
% %\Correlatore{Ubaldo Kutuzu}
% \end{frontespizio}

\begin{titlepage}
	% if you want the titlepage to be centered, uncomment and
	% fine-tune the line below (KOMA classes environment)
	% \begin{addmargin}[-1cm]{-3cm}
    \begin{figure}[!h]
    \flushleft
	\includegraphics[width=0.25\columnwidth]{images/univr.png}
	\end{figure}

    \begin{center}
    	\large
        \textsc{\large{Università degli studi di Verona}}
        \hfill
        \vfill
		\textsc{\large{Dipartimento di diagnostica e sanità pubblica}}
		\vfill
        \textsc{Dottorato di ricerca in nanoscienze e tecnologie avanzate}\\
        Curriculum in Nanomedicina Morfologica Clinica\\
		XIX Ciclo
		\vfill
        \begingroup
       		\huge\textbf{
            The modular structure of brain functional connectivity networks:\\a graph theoretical approach
            }
            \bigskip
        \endgroup
        \vfill
        S.S.D. FIS/07
		\flushleft 
		\normalsize{\textbf{Coordinatore}:}\\
		\medskip {Franco Tagliaro}
		\flushleft
		\normalsize{\textbf{Tutor}:}\\
		\medskip{Angelo Bifone}\\
		\medskip{Pasquina Marzola}\\
        \flushright
        \normalsize{\textbf{Dottorando:}}\\
        \medskip {Carlo Nicolini}\\	
		
        \vfill
        \centering{\textbf{Keywords}:\\Brain networks, modularity, community detection, functional connectivity, Surprise.}

        \vfill
		\centering {A.A. 2014-2017}

    \end{center}  
  %\end{addmargin}       
\end{titlepage}

%%%%%%%%%%%%%%%%% SECOND PAGE %%%%%%%%%%%%%%%%%
% \vspace*{\fill}
% \begin{footnotesize}
% \begin{centering}
% \noindent \copyright 2017 Carlo Nicolini \\
% \url{nicolini.carlo@gmail.com}\\
% Quest’opera è stata rilasciata con licenza Creative Commons Attribuzione – non commerciale non opere derivate 3.0 Italia. Per leggere una copia della licenza visita il sito web:

% \url{http://creativecommons.org/licenses/by-nc-nd/3.0/it/}.

% \begin{flushleft}
% \textbf{Attribuzione:} Devi riconoscere una menzione di paternità adeguata, fornire un link alla licenza ed indicare se sono state effettuate delle modifiche. Puoi fare ciò in qualsiasi maniera ragionevole possibile, ma non con modalità da suggerire che il licenziante avalli e o il tuo utilizzo del materiale.\\
% \textbf{Non commerciale:}  Non puoi usare il materiale per scopi commerciali.\\\textbf{Non opere derivate} Se remixi, trasformi il materiale o ti basi su di esso, non puoi distribuire il materiale così modificato.
% \end{flushleft}

% Graph analysis of brain functional connectivity – Carlo Nicolini \\
% Tesi di Dottorato\\
% Verona, 10 Febbraio 2017 \\
% ISBN XXX \\
% \end{centering}
% \end{footnotesize}
%%%%%%%%%%%%%%%%% SECOND PAGE %%%%%%%%%%%%%%%%%
%\thispagestyle{empty}